\section{Restricted supersingular comparisons}\label{sec:restricted}
This section proves additional comparisons for quadratic Lubin--Tate twists and for two-slope models, and determines the obstruction imposed by a stable cyclic subgroup.
\subsection{Conventions and formal-group input}
This subsection fixes the descent character and the valuation conventions used in the comparisons.


Let $K$ be a complete discretely valued field of characteristic zero with
perfect residue field of characteristic $p\geq5$. Write $e_K=v_K(p)$, where
$v_K(K^\times)=\mathbf Z$. Unramified base extension does not change this
index. When considering inertia, we pass to the completion of the maximal
unramified extension and hence assume that the residue field $k$ is
algebraically closed. Put $K_0=\operatorname{Frac} W(k)$.

A potentially supersingular Barsotti--Tate group has height two and acquires
a good dimension-one formal model with special slope $1/2$. Its special
inertia trace is the trace of finite descent on its rational covariant
Dieudonn\'e module, evaluated on prime-to-$p$ inertia. A finite torsion
isomorphism always means an isomorphism of actual generic-fibre finite
group schemes, equivalently of the corresponding finite Galois modules.

We use the height-two Lubin--Tate deformation theorem
\citep{LT1966}. For the Honda formal group, an integral universal deformation
over $W(k)[[u]]$ can be chosen in the form
\begin{equation}\label{rc-eq:ptypical}
 [p]_{F_u}(T)=pT+_{F_u}uT^p+_{F_u}T^{p^2}.
\end{equation}
The Teichm\"uller scalar action is
\begin{equation}\label{rc-eq:scalar}
 F_{a^{1-p}u}(aX,aY)=aF_u(X,Y),\qquad
 a\in\mu_{p^2-1}.
\end{equation}
This normalization and action are recorded explicitly in
\citet[\S2.1]{BDMMRR2019}. No orbit computation from that source is used.

For completeness, the tame good-reduction extension used below can be
bounded uniformly. The endomorphism algebra of the height-two simple
special isocrystal is the quaternion division algebra of invariant $1/2$.
A finite subgroup of its units has no element of order $p$: such an element
would generate a field containing a primitive $p$th root, of degree
$p-1>2$, whereas commutative subfields have degree at most two. A finite
subgroup is conjugate into the units of a maximal order. Reduction to
$\mathbf F_{p^2}^{\times}$ has pro-$p$ kernel, so the finite subgroup injects
into this cyclic group. Its order therefore divides $p^2-1$.
Killing this finite special descent action makes the potentially crystalline
representation crystalline, and the crystalline/good-reduction equivalence
for integral Barsotti--Tate representations over a complete DVR with
perfect residue field gives a good model; see
\citet[Theorem 2.2.1]{Liu2013}. Thus the tame extension of degree
$m=p^2-1$ suffices. Its special trace has the form
\begin{equation}\label{rc-eq:type}
 \theta=\alpha+\alpha^p,\qquad \alpha^m=1.
\end{equation}

\subsection{The deformation weight}

Take the tame extension $L/K$ of degree $m=p^2-1$ and write
$E=v_L(p)=me_K$. Choose a uniformizer $\pi$ of $L$ and a generator $\tau$
of $\Gal(L/K)$ with $\tau(\pi)=\zeta\pi$, where $\zeta$ is a primitive
$m$th root in $W(k)$. A special-fibre framing conjugates the finite stabilizer
into the Teichm\"uller subgroup. The representing property of the framed
deformation functor, together with \eqref{rc-eq:scalar}, then gives the
parameter action $u\mapsto\alpha^{1-p}u$.
If $\alpha(\tau)=\zeta^a$ and $u\ne0$ has valuation $h<E$, comparison of
leading residues gives
\begin{equation}\label{rc-eq:hasseweight}
 h\equiv(1-p)a\pmod{p^2-1}.
\end{equation}
In particular $p-1\mid h$. We will use this statement only below the
two-slope boundary, where the uncapped parameter valuation is well defined.

\subsection{Newton polygons and division points}

\begin{lemma}\label{rc-lem:newton}
Let $h=v_L(u)>0$, allowing $h=+\infty$. If
$h\geq pE/(p+1)$, all nonzero $p$-torsion points have valuation
$E/(p^2-1)$. If $0<h<pE/(p+1)$, put
\[
 A=\frac{E-h}{p-1},\qquad B=\frac{h}{p(p-1)}.
\]
There are $p-1$ nonzero $p$-torsion points of valuation $A$ and
$p^2-p$ of valuation $B$, with $A>B$. The points of valuation at least
$A$, together with zero, form a subgroup $C$ of order $p$.

For a nonzero torsion target of valuation $0<q\leq A<E$ in the two-slope
case, its $p$-division points have valuations
\[
 \begin{cases}
 q/p^2 & q<p^2B,\quad\text{all }p^2\text{ points};\\
 B & q=p^2B,\quad\text{all }p^2\text{ points};\\
 (q-h)/p\ \text{and}\ B & q>p^2B,
 \quad\text{respectively }p\text{ and }p^2-p\text{ points}.
 \end{cases}
\]
\end{lemma}

\begin{proof}
The lower Newton polygon of $[p](T)/T$ is determined by
$(1,E),(p,h),(p^2,0)$ before division by $T$.
Terms arising from formal addition in \eqref{rc-eq:ptypical} lie strictly
above the relevant supporting lines: a mixed monomial is a product of at
least two positive-valuation summands. The two stated polygons and root
counts follow.
The high-valuation points form a subgroup because an integral formal group
law preserves the valuation filtration. Their count is $p$ including zero.

For division of the given target, the relevant polygon has points
$(0,q),(p,h),(p^2,0)$. The linear term is above it because $q<E$.
The middle point lies above, on, or below the segment from $(0,q)$ to
$(p^2,0)$ precisely as $q<p^2B$, $q=p^2B$, or $q>p^2B$.
Reading slopes proves the three cases.
\end{proof}

The restriction $q\leq A<E$ in this lemma is essential. It is preserved
along a chain of nonzero torsion division points: for every positive
valuation $r$, the formal multiplication law gives $v([p]x)>v(x)$, so
successive divisions have strictly decreasing positive valuation.

\subsection{Unramified quadratic Lubin--Tate twists}

Let $F/\Q_p$ be unramified quadratic, let $\Ocal=\mathbf Z_{p^2}\subset W(k)$,
and let $\sigma$ be its nontrivial automorphism. Let $G_0$ be the formal
$\Ocal$-module with
\[
 [p](T)=pT+T^{p^2}.
\]
Its existence and $\Ocal$-module identities are given by
\citet[Lemma 1 and Theorem 1, pp.~381--382]{LT1965}. It is a height-two,
dimension-one supersingular Barsotti--Tate group as a $\Z_p$-group. Denote its
$\Ocal$-linear Tate character by $\psi:G_{K_0}\to\Ocal^\times$.
For a finite character $\eta:G_K\to\mu_{p^2-1}$, let $T_\eta=\Ocal$ with
character $\eta\psi|_{G_K}$. This is the standard lattice of the finite twist
of $G_0$ by $\eta$, with special trace $\eta+\eta^p$.

\begin{theorem}\label{rc-thm:LT}
Suppose
\[
 T_\eta/p^n\simeq T_{\eta'}/p^n
\]
as $\mathbf Z/p^n[G_K]$-modules. If $p^{n-1}\nmid e_K$, then
$\eta=\eta'$ with the displayed coefficient conventions, and in particular
their special inertia traces agree. This also applies to any $\Ocal$-stable
rank-one lattices in these characters.
\end{theorem}

\begin{proof}
We first derive the division-field degree needed in the argument, including
over the possibly infinite residue field $k$. Put $q=p^2$,
$f(T)=pT+T^q$, and let $f_j$ be its $j$-fold iterate. Then
\[
 Q_j(T)=\frac{f_j(T)}{f_{j-1}(T)}
       =p+f_{j-1}(T)^{q-1}
\]
is a monic Eisenstein polynomial over $W(k)$ of degree
$(q-1)q^{j-1}$. A root $P_j$ has exact order $p^j$, since
$f_{j-1}(P_j)=0$ would force $p=0$. Its annihilator in $\Ocal$ is $(p^j)$,
so its $\Ocal$-span contains $q^j$ points, all the roots of $f_j$.
The endomorphism power series $[a]$ converges in $K_0(P_j)$, so this field
contains every division point and is Galois. Its automorphisms act through
$(\Ocal/p^j)^\times$. The Eisenstein degree is the order of this unit group,
so the action is surjective. The division extension is totally ramified.

Consider the character
\[
 c_n(u)=\frac{\sigma(u)}u\quad\text{on }(\Ocal/p^n)^\times.
\]
Its image on Teichm\"uller units has order $p+1$. On principal units,
the logarithm identifies it with $\sigma-1$ on $p\Ocal/p^n\Ocal$.
Choose a unit $\delta\in\Ocal$ of trace zero; then
$\Ocal=\Z_p\oplus\Z_p\delta$ and $\sigma(\delta)=-\delta$.
Since $2$ is a unit, the image is
$p\Z_p\delta/p^n\Z_p\delta$, of order $p^{n-1}$.
Consequently the pro-$p$ part of $c_n\circ\psi$ cuts out a totally ramified
cyclic extension of $K_0$ of degree $p^{n-1}$. If that character vanishes on
$G_K$, the extension lies in $K$, forcing $p^{n-1}\mid e_K$.

Write $\Ocal_n=\Ocal/p^n$. Every $\mathbf Z/p^n$-linear map
$f:\Ocal_n\to\Ocal_n$ has a unique expression
\begin{equation}\label{rc-eq:intertwiner}
 f(x)=ax+b\sigma(x),\qquad a,b\in\Ocal_n.
\end{equation}
Indeed, evaluation at $1$ and $\delta$ gives
$a+b=f(1)$ and $a-b=f(\delta)/\delta$, determining $a,b$ uniquely.
If $f$ is invertible, at least one of $a,b$ is a unit, since otherwise
its reduction modulo $p$ is zero.

Set $\lambda=\eta\psi$ and $\lambda'=\eta'\psi$.
Intertwining and uniqueness in \eqref{rc-eq:intertwiner} give, for every $g$,
\[
 a(\lambda(g)-\lambda'(g))=0,
 \qquad b(\sigma(\lambda(g))-\lambda'(g))=0.
\]
If $a$ is a unit, $\eta\equiv\eta'\pmod p$, hence $\eta=\eta'$ because
reduction is injective on $\mu_{p^2-1}$. If $b$ is a unit, then
\[
 \frac{\sigma(\psi)}\psi=\frac{\eta'}{\sigma(\eta)}
 \pmod{p^n}\quad\text{on }G_K.
\]
The right side has prime-to-$p$ order. The pro-$p$ quotient just computed
therefore vanishes on $G_K$, contradicting $p^{n-1}\nmid e_K$.
Finally, every $\Ocal$-stable rank-one lattice is an $\Ocal$-fractional ideal
and hence isomorphic to the standard lattice.
\end{proof}

\begin{remark}
The theorem uses the full quadratic coefficient action on the lattice.
It does not assert that an arbitrary potentially supersingular group, or an
arbitrary stable $\Z_p$-lattice in such a representation, belongs to this
family. For example, it applies to the $125$-torsion comparison at
$(p,e_K)=(5,20)$ within this explicitly specified family.
\end{remark}

\subsection{A residual wild break and a second comparison theorem}\label{rc-sec:residual}
This subsection recovers a prime-to-$p$ Hasse valuation from the actual residual splitting field.


Return to the common tame field $L$ and the two-slope notation $A,B,h$
of Lemma~\ref{rc-lem:newton}.

\begin{proposition}\label{rc-prop:break}
If $p\nmid h$, the full $p$-torsion field over $L$ is cyclic of degree $p$,
with lower ramification break
\begin{equation}\label{rc-eq:break}
 b=\frac{pE-(p+1)h}{p-1}.
\end{equation}
If $p\mid h$ and the residual representation over $L$ is nontrivial, its
full splitting extension is again cyclic of degree $p$, and its break
satisfies the strict inequality
\[
 b<\frac{pE-(p+1)h}{p-1}.
\]
\end{proposition}

\begin{proof}
The canonical $p$-torsion points lie in $L$. Their valuation $A$ is integral
by \eqref{rc-eq:hasseweight}; scaling their factor by $\pi^A$ gives a nonzero
constant plus a unit times $X^{p-1}$, with simple roots in $k$.
Hensel's lemma gives all of them over $L$.
The quotient line is also constant over $L$: the determinant of a good
dimension-one Barsotti--Tate group on inertia is cyclotomic, and
$\mu_p\subset L$ because $p-1\mid E$ and $k$ is algebraically closed.
Indeed choose $b\in L$ with $b^{p-1}=-p$; after $X=bY$, the polynomial
$((1+X)^p-1)/(pX)$ reduces to $1-Y^{p-1}$, whose simple nonzero roots lift.
For the determinant background in this generality, see
\citet[proof of Proposition 5.2.2]{Liu2007}.

Choose a noncanonical point $Q$ with coordinate $x$. Galois acts by
$Q\mapsto Q+jP$ for a fixed generator $P$ of the canonical line. Thus
$L(x)$ is the full $p$-torsion field and has degree at most $p$.
If $p\nmid h$, its valuation $B=h/[p(p-1)]$ has denominator $p$, so the
degree is exactly $p$. In the second assertion nontriviality gives the same
conclusion.

Normalize $w=pv_L$ on $L(x)$. For every nontrivial $\gamma$ in its cyclic
Galois group, formal addition gives
\[
 w(\gamma(x)-x)=pA,\qquad w(x)=\frac h{p-1}.
\]
If a cyclic order-$p$ extension has lower break $b$, expansion in a
uniformizer with fixed Teichm\"uller coefficients gives
\[
 w(\gamma(y)-y)\geq w(y)+b.
\]
The inequality is an equality when $p\nmid w(y)$: the first exponent then
has nonzero derivative in the residue field. It is strict when
$p\mid w(y)$, because the first exponent's leading change vanishes and
all subsequent terms have larger valuation. Applying these statements to
$x$ proves both assertions.
\end{proof}

\begin{theorem}\label{rc-thm:primehasse}
Let $M_1,M_2$ be potentially supersingular groups over the same $K$.
Suppose $M_1[p]\simeq M_2[p]$, and over the common tame field $L$ both
models are two-slope with $p\nmid h_1h_2$. Then their special inertia traces
agree.
\end{theorem}

\begin{proof}
The actual residual representations determine the same splitting field over
$L$, hence the same lower break. Formula~\eqref{rc-eq:break} gives
$h_1=h_2=h$, and therefore the same high valuation $A$.
The canonical line is the unique fixed line of nontrivial wild inertia,
so the torsion isomorphism identifies its character $\kappa$.

For a canonical point, its coordinate is $\pi^A$ times a unit.
Field application of $\tau$ multiplies its leading residue by $\zeta^A$.
The good-model descent map has tangent multiplier $\alpha(\tau)$.
Thus the descended Galois action on the canonical line satisfies
\begin{equation}\label{rc-eq:kappa}
 \kappa(\tau)=\zeta^A\alpha(\tau)^{-1}.
\end{equation}
The same $A$ and $\kappa$ give the same $\alpha$, and
\eqref{rc-eq:type} gives equality of special traces.
\end{proof}

\begin{corollary}\label{rc-cor:strictordering}
If two compatible residual representations have two-slope parameters with
$p\nmid h_1$ and $p\mid h_2$, then $h_2<h_1$.
\end{corollary}
\begin{proof}
The common residual representation is nontrivial over $L$ by
Proposition~\ref{rc-prop:break}. Its exact formula for $h_1$ and strict upper
bound for $h_2$ imply the claimed ordering.
\end{proof}

\subsection{Stable cyclic subgroups and the threshold \texorpdfstring{$p^n$}{p\^n}}
This subsection obtains a ramification obstruction by comparing a primitive point valuation with the degree allowed by cyclic Galois action.


\begin{definition}
In the two-slope case, a cyclic subgroup $D\subset M[p^n]$ of order $p^n$
is called noncanonical at level one if its order-$p$ subgroup is one of the
low-valuation lines, rather than the canonical subgroup of
Lemma~\ref{rc-lem:newton}.
The outer canonical range is $0<h<E/(p+1)$.
\end{definition}

\begin{theorem}\label{rc-thm:cyclic}
If $D$ is $G_K$-stable and noncanonical at level one, then
$v_p(h)\geq n$. If the model is in the outer canonical range, this implies
$e_K>p^n$. Consequently no such subgroup exists there when $e_K\leq p^n$.
\end{theorem}

\begin{proof}
Let $Q$ generate $D$. Its order-$p$ multiple has valuation $B$.
Each subsequent division remains on the low branch of
Lemma~\ref{rc-lem:newton}, so
\begin{equation}\label{rc-eq:lowprimitive}
 v_L(T(Q))=\frac h{(p-1)p^{2n-1}}.
\end{equation}
The Galois action on $Q$ is faithful through a subgroup of
$(\mathbf Z/p^n)^\times$, so the $p$-part of $[L(Q):L]$ is at most
$p^{n-1}$. The $p$-part of the denominator of
\eqref{rc-eq:lowprimitive} is $p^{\max(0,2n-1-v_p(h))}$ and must divide its
ramification index. Hence $v_p(h)\geq n$.
By \eqref{rc-eq:hasseweight}, $h/(p-1)$ is a positive integer, so it is at
least $p^n$. In the outer range,
\[
 \frac h{p-1}<\frac E{(p-1)(p+1)}=e_K,
\]
which proves the strict bound.
\end{proof}

\begin{proposition}\label{rc-prop:oneslopecyclic}
If a one-slope supersingular group has a $G_K$-stable cyclic subgroup of
order $p^n$, then $v_p(e_K)\geq n-1$.
\end{proposition}
\begin{proof}
Every primitive point in this range has valuation
\[
 \frac E{(p^2-1)p^{2n-2}}=\frac{e_K}{p^{2n-2}},
\]
by repeated division along the one-slope polygon. Its denominator has $p$-part
\[p^{\max(0,2n-2-v_p(e_K))}.\] The field of a generator of a stable
cyclic subgroup has $p$-degree at most $p^{n-1}$. Comparing the two proves
the claim.
\end{proof}

These arguments do not identify the character of the stable cyclic
subgroup with a finite twist of the trivial or cyclotomic character. Such
an identification would be an additional assertion, not a consequence of
cyclicity or of the denominator estimate.

\subsection{Degree-\texorpdfstring{$p$}{p} quotient valuations}
This subsection computes how the Hasse valuation changes under each degree-$p$ quotient.


\begin{proposition}\label{rc-prop:quotients}
Let $C$ be a subgroup of order $p$ of a two-slope formal group. The Hasse
valuation $h'$ of its quotient, whenever below the one-slope boundary, is
\[
\begin{array}{c|c}
\text{case}&h'\\\hline
C\text{ canonical},\ h<E/(p+1)&ph\\
C\text{ canonical},\ h>E/(p+1)&E-h\\
C\text{ noncanonical}&h/p.
\end{array}
\]
At $h=E/(p+1)$, the canonical quotient is one-slope. A quotient of a
one-slope group by any subgroup of order $p$ has $h'=E/(p+1)$.
\end{proposition}

\begin{proof}
The quotient formal isogeny $q_C$ can be normalized, up to a unit, as the
product of $T-_{F}P$ for $P\in C$. The existence of these quotient
isogenies is classical; see \citet{Lubin1967}. For a point away from
equal-valuation cancellation this gives
\[
 v(q_C(x))=\sum_{P\in C}v(x-_{F}P).
\]

If $C$ is canonical, a low $p$-torsion point has image valuation
$pB=h/(p-1)$. A division point of a high point has valuation
$(A-h)/p$ when $A>p^2B$, and $A/p^2$ when $A<p^2B$.
Its image has valuation $p$ times that value. Since
$A>p^2B$ is equivalent to $h<E/(p+1)$, the resulting high and low
$p$-torsion valuations identify $h'$ as $ph$ in the first regime and
$E-h$ in the second. At equality all nonzero quotient $p$-torsion points
have valuation $E/(p^2-1)$.

If $C$ is noncanonical, a high point has image valuation
\[
 A+(p-1)B=\frac{E-h/p}{p-1},
\]
while division of a low point gives image valuation
$B/p=h/[p^2(p-1)]$. These identify $h'=h/p$.

Finally, in the one-slope case write $r=E/(p^2-1)$.
The distinct leading residues of the $p$-torsion form a two-dimensional
$\F_p$-space. The leading quotient polynomial has as kernel precisely the
one-dimensional residue space of $C$. Thus nonzero points of the dual
kernel have image valuation $pr$, while division points of $C$ have image
valuation $r/p$. The two quotient valuations give $h'=E/(p+1)$.
\end{proof}

In particular, a noncanonical $K$-defined subgroup of order $p$ forces
$p\mid h$, because its quotient over the same good-reduction field has an
integral Hasse valuation.

